\section*{Приложение}
\addcontentsline{toc}{section}{Приложение}
\label{sec:Apendix} \index{Apendix}

\begin{center}
    \textbf{Полный перечень гиперпараметров обучения}
\end{center}

В табл.~\ref{tab:full-hparams} приведён полный набор гиперпараметров, при которых
получены результаты раздела~\ref{sec:Chapter4}. Конфигурация вокодера HiFi-GAN
соответствует варианту V1 из~\cite{hifigan} с адаптированным под размерность
SSL-признаков входным слоём.

\begin{table}[H]
    \caption{Полный перечень гиперпараметров}
    \label{tab:full-hparams}
    \centering
    \footnotesize
    \begin{tabularx}{\textwidth}{X l}
        \toprule
        Параметр & Значение \\
        \midrule
        \multicolumn{2}{l}{\textit{Энкодер}} \\
        Скрытая размерность & 768 \\
        Частота кадров & 50 Гц (шаг 20 мс) \\
        Режим & заморожен \\
        Слой извлечения (WavLM / HuBERT) & 7 / 6 \\
        \midrule
        \multicolumn{2}{l}{\textit{Очиститель (параллельные адаптеры)}} \\
        Размерность узкого места & 256 \\
        Активация & GELU \\
        Число адаптеров & по числу слоёв энкодера \\
        Функция потерь & $L_1 + L_2$ по признакам, $\lambda=1$ \\
        \midrule
        \multicolumn{2}{l}{\textit{Вокодер}} \\
        Коэффициенты апсемплинга & [8, 8, 2, 2] \\
        Размеры ядер апсемплинга & [16, 16, 4, 4] \\
        Каналы остаточных блоков & 512 \\
        Число итераций WaveFit & 5 \\
        \midrule
        \multicolumn{2}{l}{\textit{Оптимизация}} \\
        Оптимизатор & AdamW ($\beta_1=0{,}8$, $\beta_2=0{,}99$) \\
        Начальная скорость обучения & $2\cdot10^{-4}$ \\
        Затухание скорости & экспоненциальное, $\gamma=0{,}999$ \\
        Размер батча & 16 \\
        Длительность сегмента & 2 с \\
        Шагов (адаптеры / вокодер) & 200 тыс. / 400 тыс. \\
        Аппаратура & 1 $\times$ NVIDIA RTX 3090, 24 ГБ \\
        \bottomrule
    \end{tabularx}
\end{table}
